\documentclass[11pt]{article}
\usepackage{fontspec}
\usepackage[french]{babel}
\usepackage{amsmath}
\usepackage{amssymb}
\usepackage{amsthm}
\usepackage{mathrsfs}
\usepackage{tikz-cd}
\usepackage{enumitem}
\usetikzlibrary{babel}

\newcommand{\printedpage}[1]{\marginpar{\footnotesize p.~#1}}
\newcommand{\msp}[1]{\marginpar{\footnotesize [#1]}}
\newcommand{\cercle}[1]{\tikz[baseline=(c.base)]\node[draw,circle,inner sep=1pt](c){#1};}

\begin{document}

\setcounter{footnote}{97}

\printedpage{124}
% head: A. Grothendieck : Dérivateurs, Ch. VI, Édition des Universités de Montpellier II et Paris VII
% foot: Version 9 novembre 2001

$W_S^{\mathrm{d}}$, $W_S^{\mathrm{d}}$, et le foncteur ‘oubli de $X$’ de $\underline{\mathrm{Liss}}_W X$ (même de $\mathrm{Cat}/X$) dans $\mathrm{Cat}/S$, lequel est compatible avec $W_X^{\mathrm{d}}$ et $W_S^{\mathrm{d}}$. On trouve ainsi un foncteur % [sic] W_S^d, W_S^d -> apparently W_S^d, W_X^d (or W_X^d, W_S^d), the same symbol is printed twice
\[
\left\{\begin{aligned}
&\mathrm{mult}_{X/S} : \mathrm{HOT}_W(S) \longrightarrow \mathrm{HOT}_W(S)\\
&\mathrm{mult}_{X/S}(\mathrm{hot}_{S,W}(Y)) \overset{\text{déf}}{=} \mathrm{hot}_{S,W}(X\times_S Y) \qquad \text{si $Y$ \textit{lisse} sur $S$}
\end{aligned}\right.
\]
(\footnote{Cette formule est également valable si $Y$ est dans $\underline{\mathrm{Fib}}_W\, S$, cf. $(*)$, p. 202.}), puisque $\mathrm{HOT}_W(S)$ peut s'obtenir comme localisé tant de $\underline{\mathrm{Liss}}_W S$ que de $\mathrm{Cat}/S$, pour $W_S^{\mathrm{d}}$. Quand $X$ est lui-même $W$-lisse sur $S$, ce foncteur n'est autre que
\[
\left\{\begin{aligned}
&\mathrm{mult}_\xi : \eta \longmapsto \xi\times_S \eta,\\
&\text{où }\ \xi = \mathrm{hot}_{S,W}(X).
\end{aligned}\right.
\]

Mais il n'en est pas de même si $X$ n'est pas lisse, i.e. il faut distinguer dans ce cas, même pour $Y$ $W$-lisse sur $S$,
\[
\mathrm{hot}_{S,W}(X\times_S Y) \qquad\text{et}\qquad \mathrm{hot}_{S,W}(X)\times_S \mathrm{hot}_{S,W}(Y).
\]

Le deuxième terme peut s'expliciter comme $\mathrm{hot}_{S,W}(\Phi(X)\times_S Y)$ qui reçoit $\mathrm{hot}_{S,W}(X\times_S Y)$, i.e.
\msp{page 201}
on a un homomorphisme canonique
\[
\underbrace{\mathrm{hot}_{S,W}(X\times_S Y)}_{=\,\mathrm{mult}_{X/S}(Y)} \longrightarrow \underbrace{\mathrm{hot}_{S,W}(X)\times_S \mathrm{hot}_{S,W}(Y)}_{=\,\mathrm{hot}_{S,W}(\Phi(X)\times_S Y)},
\]
induit par
\[
X\times_S Y \longrightarrow \Phi(X)\times_S Y,
\]
mais cette dernière flèche n'a aucune raison d'être dans $W_S^{\mathrm{d}}$, et n'induit donc pas nécessairement un isomorphisme dans $\mathrm{HOT}_W(S)$. (Contrexemple page 154.)

Il faut regarder à part le cas \textit{où $X$ est dans} $\underline{\mathrm{Fib}}_W S$, j'ai envie de dire que dans ce cas on a
\[
\mathrm{hot}_{S,W}(X\times_S Y) \xrightarrow{\ \sim\ } \underbrace{\mathrm{hot}_{S,W}(X)\times_S \mathrm{hot}_{S,W}(Y)}_{=\,\mathrm{hot}_{S,W}(\Phi(X)\times_S\, \Phi(Y))} \tag*{$(*)$} % ?? uncertain: under the brace "hot" is printed larger than "(Phi(X)x_S Phi(Y))", which stands at subscript level -- argument of hot_{S,W} or part of its subscript
\]
(\footnote{Cette formule n'est établie qu'en supposant ou bien $Y$ $W$-lisse (en plus de $X$ $W$-fibré sur $S$), ou bien $X$ $W$-propre, p.ex. $X$ $W$-parfait.}) pour \textit{tout} $Y$ dans $\mathrm{Cat}/S$ (pas nécessairement $W$-lisse), a fortiori que dans ce cas
\[
\left\{\begin{aligned}
&\mathrm{mult}_{X/S} = \mathrm{mult}_\xi \qquad \text{dans } \mathrm{HOT}_W(S)\\
&\text{où } \xi = \mathrm{hot}_{S/W}(X/S) \text{ (si $X$ dans $\underline{\mathrm{Fib}}_W S$).} % [sic] hot_{S/W} -> hot_{S,W}, the notation used everywhere else
\end{aligned}\right. \tag*{$(**)$}
\]

\printedpage{125}
% head: A. Grothendieck : Dérivateurs, Ch. VI, Édition des Universités de Montpellier II et Paris VII
% foot: Version 9 novembre 2001

Dans le cas où $Y$ est lisse sur $S$ (ce qui suffit pour établir $(**)$), il suffit de se rappeler que $X\longrightarrow \Phi(X)$ est dans $W_S^{\mathrm{f}}$ (résulte du fait que $X$ est $W$-fibré sur $S$ (\footnote{il suffit même $W$-prélisse sur $S$}), donc $X_s \longrightarrow s\backslash X$ est dans $W$, $\forall\ s\in S$), donc $\underbrace{X\times_S Y}_{X_Y} \longrightarrow \underbrace{\Phi(X)\times_S Y}_{\Phi(X)_Y}$ est aussi dans $W_Y^{\mathrm{f}}$, donc (comme $X_Y$, $\Phi(X)_Y$ sont $W$-fibrés sur $Y$) dans $W_Y^{\mathrm{d}}$, donc aussi dans $W_S^{\mathrm{d}}$. En fait, même sans supposer $Y$ lisse, ceci montre que l'on a
\msp{page 202}
\[
\mathrm{hot}_{S,W}(X\times_S Y) \xrightarrow{\ \sim\ } \mathrm{mult}_{Y/S}(\mathrm{hot}_{S,W}(X)) \qquad \text{si $X$ est dans $\underline{\mathrm{Fib}}_W S$.} \tag*{$(*)$}
\]

Pour avoir $(*)$ pour $Y$ quelconque, on peut par ce qui précède supposer $X$ lisse (quitte à le remplacer par $\Phi(X)$), donc $X$ dans $\underline{\mathrm{Liss}}_{0,W}S$. Il faut voir que
\[
X\times_S Y \longrightarrow X\times_S \Phi(Y)
\]
est encore dans $W_S^{\mathrm{d}}$. Or si $X$ est \textit{$W$-propre} sur $S$, le changement de base $X\longrightarrow S$ est compatible avec $W_S^{\mathrm{d}}$, $W_X^{\mathrm{d}}$ (sur $\mathrm{Cat}/S$ tout entier), et on gagne. (Ça doit être faux sans hypothèse de $W$-propreté sur $X$, \textit{ou} de $W$-lissité sur $Y$ \dots)

Quand on a
\[
f : S' \longrightarrow S,
\]
alors
\[
f^* : \mathrm{HOT}_W(S) \longrightarrow \mathrm{HOT}_W(S')
\]
est compatible avec les structures de produit.

\subsection*{\protect\cercle{8} Les foncteurs $f_!$ : théorème d'adjonction, propriétés du dérivateur $\mathrm{HOT}_W$.}

Soit
\[
f : S' \longrightarrow S
\]
dans Cat, alors le foncteur induit
\[
f_! : \mathrm{Cat}/S' \longrightarrow \mathrm{Cat}/S
\]
est compatible avec les localiseurs $(W_{S'}^{\mathrm{d}}, W_S^{\mathrm{d}})$, ou $(W_{S'}^{\mathrm{g}}, W_S^{\mathrm{g}})$, ou $(W_{S'}, W_S)$. On trouve donc un foncteur induit pour les localisées, et \texttt{[un]} diagramme commutatif
\msp{page 203}
\[
\begin{tikzcd}[column sep=small]
\mathrm{HOT}_W(S') \arrow[rrd, "\rho'"] \arrow[dd, "f_!"'] & & & & \mathrm{HOT}^{o}_W(S') \arrow[lld, "\sigma'"'] \arrow[dd, "f^{o}_!"] \\
& & \mathrm{HOT}^{\mathrm{lc}\,!}_W(S') \arrow[dd, "f^{\mathrm{lc}\,!}_!"] & & \\
\mathrm{HOT}_W(S) \arrow[rrd, "\rho"] & & & & \mathrm{HOT}^{o}_W(S). \arrow[lld, "\sigma"'] \\
& & \mathrm{HOT}^{\mathrm{lc}\,!}_W(S) & &
\end{tikzcd}
\]

\printedpage{126}
% head: A. Grothendieck : Dérivateurs, Ch. VI, Édition des Universités de Montpellier II et Paris VII
% foot: Version 9 novembre 2001

On n'a pas de foncteur $f_!^{\mathrm{lc}}$ au niveau des $\mathrm{HOT}_W^{\mathrm{lc}}$, sauf si on fait une hypothèse, telle \underline{W(4b)}, qui assure que $\mathrm{HOT}_W^{\mathrm{lc}\,!} \approx \mathrm{HOT}_W^{\mathrm{lc}}$.

Le théorème principal, que le foncteur $f_!$ est adjoint à gauche de $f^*$, est valable \textit{sous la seule hypothèse} W$(1,2,3)$ (\footnote{i.e. avec ces nouvelles notations Loc$(1,2,3\ \text{bis})$ (cf. XVI).}). Donc dualement pour $f_!^{o}$, $f_{o}^*$. Mais comme on n'a pas de foncteur $f_{\mathrm{lc}\,!}^*$, la question d'adjonction ne se pose pas pour $f$ général, et sans hypothèse W(4b) sur $W$. Elle ne se pose, apparemment, que pour $f$ un $W$-fibré trivial (cf. fin de \cercle{6}, p. 198), et dans ce cas l'adjonction ne doit pas faire de difficultés. Mais tant que je n'ai pas de sentiment bien net au sujet de la signification des $\mathrm{HOT}_W^{\mathrm{lc}\,!}$, ce n'est pas là un résultat bien bandant \texttt{[sic]}.

L'adjonction de $f_!$, $f^*$ pour $\mathrm{HOT}_W$ permet déjà de montrer que le prédérivateur $\mathrm{HOT}_W$
\msp{page 204}
satisfait les conditions Der 1, Der $4'$, Der $5'$ des dérivateurs. De plus, on a ce qu'on peut appeler la propriété Der 3b pour $\mathrm{HOT}_W$ en tant que $W$-dérivateur, i.e.

\begin{quote}
\textit{Si $f\in W$, alors $f$ est une} $\mathrm{HOT}_W$\textit{-équivalence} (\footnote{On le prouve en interprétant la notion de $\mathrm{HOT}_W$-équivalence par la $\mathrm{HOT}_W$\textit{-homologie.} Cela revient finalement à ceci : Si $f : S' \longrightarrow S$ est dans $W$, alors $\forall\, T$, $S'\times T \longrightarrow S\times T$ aussi!}) (moyennant les seuls axiomes W$(1,2,3)$ sur $W$). Plus précisément, $f$ est une $\mathrm{HOT}_W$-équivalence si et seulement si $\mathrm{hot}_W(f)$ est \texttt{[un]} isomorphisme, i.e. $f\in \overline{W}$, le saturé fort de $W$.
\end{quote}

Si $W$ \textit{satisfait} W(4b)\textit{,} alors on a également un couple de foncteurs adjoints
\[
\begin{tikzcd}
\mathrm{HOT}_W^{\mathrm{lc}}(S') \arrow[r, shift left, "f_!^{\mathrm{lc}}"] & \mathrm{HOT}_W^{\mathrm{lc}}(S) \arrow[l, shift left, "f^*_{\mathrm{lc}}"]
\end{tikzcd}
\]
pout toute flèche $S'\longrightarrow S$ dans Cat. On voit aussi dans ce cas que % [sic] pout -> pour, misprint

\begin{quote}
% printed with a vertical rule at the left of the block
\begin{tabular}{|p{0.9\linewidth}}
si $f : S' \longrightarrow S$ est dans $W$, alors $f^*_{\mathrm{lc}}$ et $f_!^{\mathrm{lc}}$ entre $\mathrm{HOT}_W^{\mathrm{lc}}(S)$ et $\mathrm{HOT}_W^{\mathrm{lc}}(S')$ sont des équivalences quasi-inverses l'une de l'autre (moyennant W(4b)),
\end{tabular}
\end{quote}

ce qui est l'axiome Der 3b$_2$ ‘version $W$-dérivateurs’. On ne peut s'attendre à mieux, du côté Der 3!

\textbf{Oubli :} Soit $f : S' \longrightarrow S$, on a alors $\boxed{f_!f^*} : \mathrm{HOT}_W(S) \longrightarrow \mathrm{HOT}_W(S)$ est égal à $\mathrm{mult}_{S'/S}$ (cf. p. 200). Il est donc égal à $\mathrm{mult}_\xi$, où $\xi = \mathrm{hot}_{S,W}(S')$, si on suppose ou bien $S'$ $W$-lisse sur $S$, ou bien $S'$ $W$-fibré sur $S$ (cf. $(**)$, p. 201).

\subsection*{\protect\cercle{9} Théorème de changement de base pour $f_!$ dans $\mathrm{HOT}_W$.}

Si on a
\printedpage{127}
% head: A. Grothendieck : Dérivateurs, Ch. VI, Édition des Universités de Montpellier II et Paris VII
% foot: Version 9 novembre 2001
\[
\begin{tikzcd}
X \arrow[d, "f"'] & X' \arrow[l, "q"'] \arrow[d, "f'"] \\
S & S', \arrow[l, "g"']
\end{tikzcd}
\]
avec $f$ $W$-lisse \textit{ou} $g$ $W$-propre \texttt{[la suite manque]}

\end{document}
